% GENERATED FILE. Edit canonical lessons, then run npm run generate:books.
% canonical-source-hash: fnv1a64:07ef1db722b054d8
% canonical-lessons: TA-C479-tacai, TA-C479-narampu, TA-C479-nencu, TA-C479-vila, TA-C479-tuti

\chapter{Muscles and Nerves}
\label{ch:ta-muscles-and-nerves}
\hlchaptermodality{479}

\begin{chapteropening}
\textbf{By the end of this chapter:} \emph{I can talk about muscles, nerves, the chest and the ribs.}

Complete the last lesson of chapter 479: I can talk about muscles, nerves, the chest and the ribs.
\end{chapteropening}

\section[\texorpdfstring{\ta{தசை} (tacai)}{தசை (tacai)}]{\ta{தசை} (tacai) — a muscle; flesh}

\label{lesson:TA-C479-tacai}



\begin{quote}
Before the new one: say the Tamil for a kidney, then the Tamil for the intestines.
\end{quote}

\subsection*{You'll want to know: \ta{தசை}}
\textbf{\ta{தசை}} — \emph{tacai} — \textquotedblleft{}a muscle; flesh\textquotedblright{}.

\begin{quote}
\ta{உடற்பயிற்சி} \ta{செய்தால்} \ta{தசைகள்} \ta{உறுதியாகும்}.
\end{quote}

\begin{quote}
\emph{uṭaṟpayiṟci ceytāl tacaikaḷ uṟutiyākum.}
\end{quote}

\begin{quote}
If you exercise, your muscles get firm.
\end{quote}

Plural: \ta{தசைகள்}, \emph{tacaikaḷ}.

Plural \ta{தசைகள்}.

\begin{cousinweb}
\textbf{\ta{தசை}} (\emph{tacai}) is native Dravidian (DEDR 3016), for flesh, muscle and the pulp of fruit.
\end{cousinweb}

\begin{grammarlens}[title={What you've built}]
One more word for describing what you lived and felt.
\end{grammarlens}

\subsection*{Your turn}
\begin{itemize}
\raggedright
  \item \emph{Say it:} \emph{tacai}
  \item \emph{Say it:} \emph{tacai}, once more
  \item \emph{Say it:} \emph{uṭaṟpayiṟci ceytāl tacaikaḷ uṟutiyākum.} — If you exercise, your muscles get firm.
  \item \emph{Recall:} say the Tamil for a kidney, then the Tamil for the intestines, then say \emph{tacai} again
\end{itemize}

\begin{tcolorbox}[breakable,colback=teal!4,colframe=teal!35!black,title={Before you move on}]
What does \ta{தசை} mean? (\textquotedblleft{}A muscle; flesh\textquotedblright{}.) Say it once more.
\end{tcolorbox}
\section[\texorpdfstring{\ta{நரம்பு} (narampu)}{நரம்பு (narampu)}]{\ta{நரம்பு} (narampu) — a nerve; a vein}

\label{lesson:TA-C479-narampu}



\begin{quote}
Before the new one: say the Tamil for the intestines, then the Tamil for a muscle.
\end{quote}

\subsection*{You'll want to know: \ta{நரம்பு}}
\textbf{\ta{நரம்பு}} — \emph{narampu} — \textquotedblleft{}a nerve; a vein\textquotedblright{}.

\begin{quote}
\ta{கையில்} \ta{உள்ள} \ta{நரம்பில்} \ta{ஊசி} \ta{போட்டார்கள்}.
\end{quote}

\begin{quote}
\emph{kaiyil uḷḷa narampil ūci pōṭṭārkaḷ.}
\end{quote}

\begin{quote}
They gave the injection into a vein in the arm.
\end{quote}

One word for nerve, sinew and vein.

\begin{cousinweb}
\textbf{\ta{நரம்பு}} (\emph{narampu}) is native Dravidian (DEDR 2903), covering nerve, tendon and blood vessel.
\end{cousinweb}

\begin{grammarlens}[title={What you've built}]
One more word for describing what you lived and felt.
\end{grammarlens}

\subsection*{Your turn}
\begin{itemize}
\raggedright
  \item \emph{Say it:} \emph{narampu}
  \item \emph{Say it:} \emph{narampu}, once more
  \item \emph{Say it:} \emph{kaiyil uḷḷa narampil ūci pōṭṭārkaḷ.} — They gave the injection into a vein in the arm.
  \item \emph{Recall:} say the Tamil for the intestines, then the Tamil for a muscle, then say \emph{narampu} again
\end{itemize}

\begin{tcolorbox}[breakable,colback=teal!4,colframe=teal!35!black,title={Before you move on}]
What does \ta{நரம்பு} mean? (\textquotedblleft{}A nerve; a vein\textquotedblright{}.) Say it once more.
\end{tcolorbox}
\section[\texorpdfstring{\ta{நெஞ்சு} (neñcu)}{நெஞ்சு (neñcu)}]{\ta{நெஞ்சு} (neñcu) — the bosom; the chest; the heart}

\label{lesson:TA-C479-nencu}



\begin{quote}
Before the new one: say the Tamil for a muscle, then the Tamil for a nerve.
\end{quote}

\subsection*{You'll want to know: \ta{நெஞ்சு}}
\textbf{\ta{நெஞ்சு}} — \emph{neñcu} — \textquotedblleft{}the bosom; the chest; the heart\textquotedblright{}.

\begin{quote}
\ta{படிக்கட்டு} \ta{ஏறியதும்} \ta{எனக்கு} \ta{நெஞ்சு} \ta{வலித்தது}.
\end{quote}

\begin{quote}
\emph{paṭikkaṭṭu ēṟiyatum eṉakku neñcu valittatu.}
\end{quote}

\begin{quote}
After climbing the stairs I had a pain in my chest.
\end{quote}

Also the heart as the seat of feeling: \ta{நெஞ்சு} \ta{நிறைந்த} \ta{நன்றி}, heartfelt thanks.

\begin{cousinweb}
\textbf{\ta{நெஞ்சு}} (\emph{neñcu}) is native Dravidian (DEDR 3736), for the chest and the heart as the seat of courage and feeling.
\end{cousinweb}

\begin{grammarlens}[title={What you've built}]
One more word for describing what you lived and felt.
\end{grammarlens}

\subsection*{Your turn}
\begin{itemize}
\raggedright
  \item \emph{Say it:} \emph{neñcu}
  \item \emph{Say it:} \emph{neñcu}, once more
  \item \emph{Say it:} \emph{paṭikkaṭṭu ēṟiyatum eṉakku neñcu valittatu.} — After climbing the stairs I had a pain in my chest.
  \item \emph{Recall:} say the Tamil for a muscle, then the Tamil for a nerve, then say \emph{neñcu} again
\end{itemize}

\begin{tcolorbox}[breakable,colback=teal!4,colframe=teal!35!black,title={Before you move on}]
What does \ta{நெஞ்சு} mean? (\textquotedblleft{}The bosom; the chest; the heart\textquotedblright{}.) Say it once more.
\end{tcolorbox}
\section[\texorpdfstring{\ta{விலா} (vilā)}{விலா (vilā)}]{\ta{விலா} (vilā) — a rib; the side of the body}

\label{lesson:TA-C479-vila}



\begin{quote}
Before the new one: say the Tamil for a nerve, then the Tamil for the bosom.
\end{quote}

\subsection*{You'll want to know: \ta{விலா}}
\textbf{\ta{விலா}} — \emph{vilā} — \textquotedblleft{}a rib; the side of the body\textquotedblright{}.

\begin{quote}
\ta{சிரித்துச்} \ta{சிரித்து} \ta{எனக்கு} \ta{விலா} \ta{வலித்தது}.
\end{quote}

\begin{quote}
\emph{cirittuc cirittu eṉakku vilā valittatu.}
\end{quote}

\begin{quote}
I laughed so much that my sides hurt.
\end{quote}

\ta{விலா} \ta{எலும்பு}, a rib bone.

\begin{cousinweb}
\textbf{\ta{விலா}} (\emph{vilā}) is native Dravidian (DEDR 5425); Malayalam has the same word for the side and the chest.
\end{cousinweb}

\begin{grammarlens}[title={What you've built}]
One more word for describing what you lived and felt.
\end{grammarlens}

\subsection*{Your turn}
\begin{itemize}
\raggedright
  \item \emph{Say it:} \emph{vilā}
  \item \emph{Say it:} \emph{vilā}, once more
  \item \emph{Say it:} \emph{cirittuc cirittu eṉakku vilā valittatu.} — I laughed so much that my sides hurt.
  \item \emph{Recall:} say the Tamil for a nerve, then the Tamil for the bosom, then say \emph{vilā} again
\end{itemize}

\begin{tcolorbox}[breakable,colback=teal!4,colframe=teal!35!black,title={Before you move on}]
What does \ta{விலா} mean? (\textquotedblleft{}A rib; the side of the body\textquotedblright{}.) Say it once more.
\end{tcolorbox}
\section[\texorpdfstring{\ta{துடி} (tuṭi)}{துடி (tuṭi)}]{\ta{துடி} (tuṭi) — to throb; to writhe}

\label{lesson:TA-C479-tuti}



\begin{quote}
Before the new one: say the Tamil for the bosom, then the Tamil for a rib.
\end{quote}

\subsection*{You'll want to know: \ta{துடி}}
\textbf{\ta{துடி}} — \emph{tuṭi} — \textquotedblleft{}to throb; to writhe\textquotedblright{}.

\begin{quote}
\ta{வலியில்} \ta{அவன்} \ta{துடித்தான்}.
\end{quote}

\begin{quote}
\emph{valiyil avaṉ tuṭittāṉ.}
\end{quote}

\begin{quote}
He writhed in pain.
\end{quote}

Past (I …): \ta{துடித்தேன்}, \emph{tuṭittēṉ}.

Also of the heartbeat (\ta{இதயம்} \ta{துடிக்கிறது}) and of eagerness: \ta{போகத்} \ta{துடித்தான்}, he was itching to go.

\begin{cousinweb}
\textbf{\ta{துடி}} (\emph{tuṭi}) is native Dravidian (DEDR 3294). DEDR keeps \ta{துடி} \textquotedblleft{}lip\textquotedblright{} and \ta{துடி} \textquotedblleft{}a small drum\textquotedblright{} apart as other entries.
\end{cousinweb}

\begin{grammarlens}[title={What you've built}]
One more word for describing what you lived and felt.
\end{grammarlens}

\subsection*{Your turn}
\begin{itemize}
\raggedright
  \item \emph{Say it:} \emph{tuṭi}
  \item \emph{Say it:} \emph{tuṭi}, once more
  \item \emph{Say it:} \emph{valiyil avaṉ tuṭittāṉ.} — He writhed in pain.
  \item \emph{Recall:} say the Tamil for the bosom, then the Tamil for a rib, then say \emph{tuṭi} again
\end{itemize}

\begin{tcolorbox}[breakable,colback=teal!4,colframe=teal!35!black,title={Before you move on}]
What does \ta{துடி} mean? (\textquotedblleft{}To throb; to writhe\textquotedblright{}.) Say it once more.
\end{tcolorbox}
